% Continuation of the Legendre audit, reserved IDs A.34--A.37.
\subsection{A pairing calculation at a multiplicative place}

\begin{statement}{A.34}{Proved}{A full integral Tate comparison fixes the multiplier}
Let $F$ be a finite extension of $\Q_p$, and let $E_1,E_2/F$ have
potentially multiplicative reduction.  Suppose an isomorphism
\[
                g:\prod_\ell T_\ell(E_1)
                     \longrightarrow\prod_\ell T_\ell(E_2)
\]
of continuous $\widehat\Z[G_F]$-modules is given.  Its Weil multiplier
is a square in $\widehat\Z^\times$.
Consequently, for a section in A.14, the multiplier obstruction A.17
vanishes whenever a potentially multiplicative place of its realized
Legendre curve is an ordinary point place of the section.
\end{statement}
\begin{proof}
Take a common finite extension $F'/F$ over which both curves are split
Tate curves, with parameters $q_1,q_2\in F'^\times$ and positive
integer normalized valuations $m_i=v_{F'}(q_i)>0$.  The curve
isomorphisms over $F'$ preserve the canonical principal polarizations,
so restriction and these changes of model do not change the multiplier.
For every coefficient prime, including $\ell=p$, the Tate sequence is
\[
       0\longrightarrow\Z_\ell(1)\longrightarrow
        T_\ell(E_i)\longrightarrow\Z_\ell\longrightarrow0,
\tag{A.34.1}
\]
with extension class $\kappa_\ell(q_i)$.  Its image under the Kummer
valuation map to $\Z_\ell$ is $m_i\ne0$, so its rational extension
is nonsplit.  The cyclotomic line is therefore the unique invariant
line of cyclotomic character: another such line or a complementary
trivial line would split the extension.  The two characters are
distinct on $G_{F'}$.  This argument also treats the coefficient prime
$p$ directly and does not discard its extension class.

The line in (A.34.1) is saturated.  In the Tate uniformization the
torsion representatives satisfy the actual pairing formula
\[
 e_N(\zeta_N^a q_i^{b/N},\zeta_N^c q_i^{d/N})
                         =\zeta_N^{ad-bc}.
\]
This is the universal Tate-curve identity and specializes to each
$F'$-curve (with one common sign convention for the Weil pairing);
see \cite[\S19.2]{LegIovitaMasdeu2009}.
Thus the Tate bases used for the extension classes are also compatible
with the Weil pairings, and the restriction of $g_\ell$ has matrix
\[
                \begin{pmatrix}a_\ell&b_\ell\\0&d_\ell\end{pmatrix},
                    \qquad a_\ell,d_\ell\in\Z_\ell^\times.
\]
Comparing the upper-right cocycles gives
\[
       \kappa_\ell(q_2)=(a_\ell/d_\ell)\kappa_\ell(q_1),
       \qquad m_2=(a_\ell/d_\ell)m_1.
\tag{A.34.2}
\]
In particular $m_2/m_1$ is a unit at every rational prime $\ell$.
It is a positive rational number, so it equals one.  Equation
(A.34.2) then gives $a_\ell=d_\ell$ for every $\ell$.  The multiplier
is $\det g_\ell=a_\ell^2$, proving the assertion.  Finite quadratic
splitting twists were removed only by their actual isomorphisms over
$F'$; a root-of-unity factor or a unit in a Tate parameter does not
alter the valuation calculation.  The full extension comparison,
rather than merely its semisimplification, is used in (A.34.2).

For the global application, an ordinary local section at the chosen
prime is the full integral Tate system of $E_{x_p}/\Q_p$.  The
isomorphism A.16 identifies it with that of $E_{\lambda_0}$.
The latter's infinite unipotent inertia remains infinite under this
isomorphism, so $E_{x_p}$ is also potentially multiplicative.
Apply the local assertion and then A.17.
\end{proof}

\begin{statement}{A.35}{Proved}{The exact remaining local locus}
A rational Legendre curve with no potentially multiplicative odd
prime has parameter $\lambda\in\{-1,\tfrac12,2\}$ and $j=1728$.
For every other realization in A.16, a multiplier class not yet
settled by A.34 can occur only when the section is locally cuspidal
at every potentially multiplicative prime of the realized curve.
\end{statement}
\begin{proof}
At an odd prime, (A.14.1) has negative valuation whenever
$v_p(\lambda)>0$, $v_p(1-\lambda)>0$, or $v_p(\lambda)<0$.
Such a curve is potentially multiplicative.  The stated absence
therefore forces both $\lambda$ and $1-\lambda$ to belong to
$\{\pm2^n:n\in\Z\}$.  Write them as $\epsilon2^a,\eta2^b$.
If $a=b$, their sum one forces $a=b=-1$ and both signs positive.
If $a\ne b$, valuation of the sum forces the smaller exponent to be
zero; the resulting equation is $-1+2=1$.  This gives exactly
$\lambda=-1,1/2,2$.  Substitution in (A.14.1) gives $j=1728$.
Every other realized parameter has a potentially multiplicative odd
place, and A.34 applies at any such place which is ordinary for the
section.  At an exact cusp the contact valuation is an element of
$\widehat\Z$, rather than the positive integer $m_2$ in A.34, so the
positive-rational-unit argument has not been supplied there.
\end{proof}

\subsection{The Gaussian invariant-line case, including the CM remainder}

\begin{statement}{A.36}{Proved}{An invariant line over $\Q(i)$ recognizes a section}
Let $s$ be a Selmer section of $U/\Q$.  If
$\bar\rho_{s,5}|_{G_{\Q(i)}}$ has an invariant line, then $s$ is a
rational point or cuspidal section.  In particular, every noncuspidal
section in A.14 whose realization has $j=1728$ is a rational point
section.  Thus that case of A.35 is completely recognized.
\end{statement}
\begin{proof}
We first verify the arithmetic of $X_0(20)$ over $K=\Q(i)$, without
an analytic-rank table.  Its genus is one by the calculation in A.13,
and its rational cusp gives its elliptic-curve origin.  The space of
weight-two cusp forms of level 20 is one-dimensional.  Its normalized
form and the form of its quadratic $-1$ twist are
\[
 f_{20}(\tau)=\eta(2\tau)^2\eta(10\tau)^2,\qquad
 f_{80}(\tau)=
 \frac{\eta(4\tau)^6\eta(20\tau)^6}
 {\eta(2\tau)^2\eta(8\tau)^2\eta(10\tau)^2\eta(40\tau)^2}.
\tag{A.36.1}
\]
The opened primary table in Kaneko--Sakai gives the associated models
\[
 A:y^2=x^3+x^2-36x-140,
 \qquad A^{(-1)}:y^2=x^3-x^2-36x+140,
\]
of conductors 20 and 80 respectively
\cite[main theorem and second table, printed p.~5]{LegKanekoSakai2013}.
The dimension-one assertion and Eichler--Shimura put $X_0(20)$ in
the isogeny class of $A$.  Both functions in (A.36.1) are strictly
positive on $\tau=iy$, since every eta factor there is positive.
They are cusp forms, so the Mellin integrals converge at both ends.
Consequently
\[
 L(A,1)=2\pi\int_0^\infty f_{20}(iy)\,dy>0,
 \qquad
 L(A^{(-1)},1)=2\pi\int_0^\infty f_{80}(iy)\,dy>0.
\tag{A.36.2}
\]
This proves analytic rank zero for both isogeny classes.

The field change used next is explicit.  Write
$B=\operatorname{Res}_{K/\Q}(A_K)$ and let
$\theta:A^{(-1)}_K\to A_K$, $(x,y)\mapsto(-x,iy)$, be the twist isomorphism satisfying
$\sigma\theta=-\theta$ for the nontrivial $\sigma\in\operatorname{Gal}(K/\Q)$.
Over $K$, the map
\[
       A\times A^{(-1)}\longrightarrow B,
       \qquad(P,Q)\longmapsto(P+\theta Q,P-\theta Q)
\tag{A.36.3}
\]
is equivariant for descent, since $\sigma$ swaps the two factors of
$B_K\simeq A_K\times A_K$ and changes the sign of $\theta$.
It descends to a $\Q$-isogeny.  The reverse map over $K$ is
$(R,S)\mapsto(R+S,\theta^{-1}(R-S))$; both compositions are
multiplication by two.  Equations (A.36.2), rank-zero modularity,
and finite Selmer groups imply finite Mordell--Weil and
Tate--Shafarevich groups for $B$.  Finiteness of the latter is
unchanged by these isogenies: the kernel of the map induced by either
reverse isogeny is contained in the finite 2-torsion of the relevant
Tate--Shafarevich group.  Shapiro's maps on global and local cohomology
identify
\[
                 B(\Q)=A(K),\qquad
                 \operatorname{Sha}(\Q,B)
                       =\operatorname{Sha}(K,A_K).
\]
The same finiteness holds for $X_0(20)_K$ by its isogeny with $A_K$.
Stoll's Corollary 6.2 and Theorem 8.6 now give the finite-descent
property over $K$ \cite{LegStoll2007,LegStollErrata2017}.

Restrict the section to $G_K$.  Its invariant line gives the actual
lift to $Z_{5,K}$ and the finite nonconstant morphism
$Z_{5,K}\to X_0(20)_K$ constructed in (A.13.1).
The completed local centres map to one point of $X_0(20)(K)$.
Their inverse image is finite, so finite-support recognition over
$K$ makes the lifted section a point or cusp section.  Thus $s_K$
has the same property.

We also spell out descent of that conclusion.  If $s_K$ is the section
of $x\in U(K)$, then its class is invariant under
$\operatorname{Gal}(K/\Q)$ because it is the restriction of $s$.
Applying the $t$-Kummer coordinate gives
$\kappa_K(x)=\kappa_K(\sigma x)$.
The map $K^\times\to\widehat{K^\times}$ is injective: valuations
first force an element in its kernel to be a unit, and the finitely
generated unit group then has trivial intersection of all power
subgroups.  Hence $x\in\Q$.
At every finite prime, pass to a completion of $K$.  If the local
restriction of $s$ is an interior point, local Kummer injectivity
identifies it with $x$.  If it is a cusp, use the unit $1-t$, $t$, or
$t/(t-1)$, respectively: it has residue one at that cusp and never
takes the value one on $U$, so it cannot restrict to the point section
of $x$.  All finite local centres of $s$ are therefore $x$, and
finite-support recognition gives $s=s_x$.
If $s_K$ is cuspidal, the same three units exclude every interior
local centre downstairs.  Its finite local centres are all supported
on $\{0,1,\infty\}$, so finite-support recognition makes $s$
cuspidal.  This proves the invariant-line assertion without an
unstated descent theorem for sections.

Finally, for each $\lambda_0=-1,1/2,2$, translation to the centre of
the three roots gives a model $y^2=x^3-dx$ over $\Q$.
The automorphism $(x,y)\mapsto(-x,iy)$ is defined over $K$.
Since 5 splits in $\Z[i]$, choose $\iota\in\Z_5$ with
$\iota^2=-1$.  The integral idempotents
$(1\pm[i]/\iota)/2$ split $T_5(E_{\lambda_0})$ into
$G_K$-stable rank-one summands.  Transport through A.16 and reduction
modulo five give the invariant line required above.  This establishes
the claimed CM case.
\end{proof}
The only multiplier cases left by A.34--A.36 are non-CM realizations
for which every potentially multiplicative place is locally cuspidal.
Even in the ordinary-place cases where A.34 proves an $H$-conjugator,
recognition from the full congruence-kernel quotient remains the
separate open step in A.24.

\begin{statement}{A.37}{Proved}{One cusp label and one bad place are recognized}
A noncuspidal Selmer section in A.14 whose locally cuspidal centres
have at most one of the three labels is a rational point section.
Every locally cuspidal place is a potentially multiplicative place
of the realized elliptic curve.  In particular, every such section
whose realization has at most one potentially multiplicative prime
is a rational point section.
\end{statement}
\begin{proof}
Let $c$ be the single possible label, choosing any label if there are
no cuspidal places.  Choose respectively
\[
       f=1-t\quad(c=0),\qquad f=t\quad(c=1),\qquad
       f=t/(t-1)\quad(c=\infty).
\]
This is a rational degree-one function, a unit on $U$, regular at
$c$ with value one.  Its pullback Kummer class is represented by an
actual nonzero local number everywhere: by $f(x_p)$ at an interior
point and by one at a cusp.  A.14 makes these numbers units outside
a finite set.  A.10 therefore makes the class one rational number
$a\in\Q^\times$.  If any local cusp occurs, injectivity at that
prime gives $a=1$.  Then no ordinary local point can occur, since
none of these three functions takes the value one on $U$; this
contradicts density-one ordinary integrality.  Thus there are no
local cusps.  All local points equal the rational point $f^{-1}(a)$,
which is interior, and finite-support recognition gives its point
section.  This is also the single-label argument in AP-D.31, with
its hypotheses now supplied by A.14.

For the second assertion, suppose the realization were potentially
good at a locally cuspidal prime $p$.  Choose an auxiliary coefficient
$\ell\ne p$ and a common finite local extension over which the
elliptic curve has good reduction and the cusp's quadratic character
is trivial.  The cusp semisimplification has Frobenius trace $1+q$,
where $q>1$ is the residue-field cardinality.  The elliptic curve has
trace of absolute value at most $2\sqrt q$, which is strictly smaller.
This contradicts the local restriction of the global Tate-system
isomorphism.  Hence the curve is potentially multiplicative there.
With at most one such prime, there is at most one local cusp and
therefore at most one label, reducing to the first assertion.
\end{proof}
For example the realized positive integral parameters $2,3,4,5$ are
recognized by A.36 or A.37: the potentially multiplicative supports
are respectively $\varnothing,\{3\},\{3\},\{5\}$.
The first positive integral parameter whose support has two primes is
$\lambda_0=6$, with support $\{3,5\}$.
This identifies a small remaining realization locus; it asserts
neither existence of a nonpoint section there nor failure of recognition.
